\documentclass[../main.tex]{subfiles}
\zexternaldocument*{category_theory}

\begin{document}

\chapter{Sheaf theory}

\section{Introduction}

\todo[inline, color=blue!40]{Write motivations for sheafs and presheafs}

\section{Presheaves and sheaves}
\subsection{Presheaves}
\begin{definition}[Presheaf] \label{def:presheaf}
  Let \(X\) be a topological space and \(\open(X)\) be its associated \nameref{ex:open-set-category}.
  For an arbitray \nameref{def:category} \(\mathsf{C}\), a \(\mathsf{C}\)-\textbf{valued presheaf} over \(X\) is a \nameref{def:functor}
  \begin{equation*}
    \mathcal{F}: \open(X)^{\op} \to \mathsf{C}
  \end{equation*}
  Less abstractly, \(\mathcal{F}\) consists of:
  \begin{enumerate}
    \item For every open set \(U \subseteq X\) an object \(\mathcal{F}(U)\) in \(\mathsf{C}\), named \textit{sections} of \(\mathcal{F}\) over \(U\).
    \item For every inclussion of open sets \(V \subseteq U\) of \(X\), a morphism \(\res_{U, V}: \mathcal{F}(U) \to \mathcal{F}(V)\) in \(\mathsf{C}\) named \textit{restriction}, satisfying \(\res_{U, U} = \id_{U}\) and for every \(W \subseteq V \subseteq U\) open subsets \(\res_{V, W} \circ \res_{U, V} = \res_{U, V}\).

    If the source object is clear, we will write \(s|_{V}\) for \(\res_{U, V}(s)\).
  \end{enumerate}
\end{definition}
\SetPluralName{def:presheaf}{Presheaves}

\begin{example}[Constant \nameref{def:presheaf}] \label{ex:constant-presheaf}
  For an arbitrary set \(S\), the \textbf{constant presheaf} (of sets) associated to \(S\) is the \nameref{def:presheaf} defined as \(\underline{S}_{\pre}(U) = S\) for all \(U \in \open(X)\).

  Note that this verifies the presheaf conditions trivially since for every open set \(U \in \open(X)\), \(\underline{S}_{\pre}(U)\) is a set (namely \(S\)) and, if \(V \subseteq U\) then we can pick \(\id_{S}\) as \(\res_{U, V}\).
\end{example}
\SetPluralName{ex:constant-presheaf}{Constant presheaves}

\subsection{Sheaves}
\begin{definition}[Sheaf] \label{def:sheaf}
  Let \(X\) be a topological space.
  A \nameref{def:presheaf} of sets \(\sheaf{F}\) is a \textbf{sheaf} if for every open set \(U \subseteq X\) and open cover \(\{U_{i}\}_{i \in \mathcal{I}}\) of \(U\) the following axioms hold:
  \begin{enumerate}[label=SH.\arabic*), ref=SH.\arabic*]
    \item\label{axiom:sheaf-identity} \textbf{Identity}: Given two sections \(s, t \in \sheaf{F}(U)\) such that for every \(i \in \mathcal{I} \quad s|_{U_{i}} = t|_{U_{i}}\) we have that \(s = t\).
    \item\label{axiom:sheaf-gluability} \textbf{Gluability}: Given a collection of sections \(\{s_{i} \mid s_{i} \in \sheaf{F}(U_{i})\}_{i \in \mathcal{I}}\) that agree on pairwise intersections (\(\forall i, j \in \mathcal{I} \;\; s_{i}|_{U_{i} \cap U_{j}} = s_{j}|_{U_{i} \cap U_{j}}\)) there exists a section \(s \in \sheaf{F}(U)\) such that \(\forall i \in \mathcal{I} \;\; s|_{U_{i}} = s_{i}\).
  \end{enumerate}
\end{definition}
\SetPluralName{def:sheaf}{Sheaves}

\begin{notation}
  To differentiate between a \nameref{def:presheaf} and a \nameref{def:sheaf}, we will denote the former using \(\mathcal{F}, \mathcal{G}, \mathcal{H}, \ldots\) and the later using \(\sheaf{F}, \sheaf{G}, \sheaf{H}, \ldots\)
\end{notation}

We will now see an example of a \nameref{def:presheaf} that is not a \nameref{def:sheaf}.
\begin{example}[The \nameref{ex:constant-presheaf} is not a \nameref{def:sheaf}] \label{ex:constant-presheaf-is-not-sheaf}
  Consider a set \(X = \{\triangle, \square\}\) endowed with the discrete topology and let \(\underline{S}_{\pre}\) be the \nameref{ex:constant-presheaf} over \(X\) of an arbitrary set \(S\) with more than one element, say \(S = \{0, 1\}\) with \(0 \neq 1\).

  Consider the open cover \(\{X_{1} = \{\triangle\}, X_{2} = \{\square\}\}\) of \(X\) and define \(0 = s_{1} \in \underline{S}_{\pre}(X_{1}) = S\) and \(1 = s_{2} \in \underline{S}_{\pre}(X_{2}) = S\).
  Since \(X_{1} \cap X_{2} = \varnothing\), \(s_{1}|_{X_{1} \cap X_{2}} = s_{2}|_{X_{1} \cap X_{2}}\).

  Now, let \(s \in \underline{S}_{\pre}(X)\) such that \(s|_{U_{1}} = s_{1}\) and \(s|_{U_{2}} = s_{2}\).
  Since the restrictions from \(\underline{S}_{\pre}(X)\) are the identity maps we have that
  \begin{equation*}
    0 = s_{1} = s|_{U_{1}} = s = s|_{U_{2}} = s_{2} = 1
  \end{equation*}
  which is a contradiction because we assumed \(0 \neq 1\).

  Therefore, since \ref{axiom:sheaf-gluability} fails, \(\underline{S}_{\pre}\) is not a sheaf.
\end{example}


However, we can make it a sheaf by considering \textit{locally constant} applications instead.
\begin{example}[Constant sheaf] \label{ex:constant-sheaf}
  Let \(X\) be a topological space and \(A\) be an object of a \nameref{def:category} \(\mathsf{C}\)\unsure{What type of category?}.
  For every open set \(U \subseteq X\), define the \textbf{constant sheaf} over \(X\) of \(A\) as
  \begin{equation*}
    \underline{S}(U) = \{f: U \to A \mid f \text{ locally constant}\}
  \end{equation*}
  Note that this is a \nameref{def:presheaf}.

  To verify that this is indeed a \nameref{def:sheaf}, let \(U \subseteq X\) be a non-empty open set an \(\{U_{i}\}_{i \in \mathcal{I}}\) be an open cover of \(U\) (in the case where \(U = \varnothing\) the axioms hold trivially).
  \begin{itemize}
    \item[\ref{axiom:sheaf-identity}] Suppose \(s\) and \(t\) are two sections of \(\underline{S}(U)\) such that \(\forall i \in \mathcal{I} \;\; s|_{U_{i}} = t|_{U_{i}}\).

          Note that for every \(x \in U\) there exists \(U_{i}\) such that \(x \in U_{i}\).

          Therefore \(s|_{U_{i}} = \{x\} = t_{U_{i}}\) and, since this holds for every \(x \in U\), \(s = t\).

    \item[\ref{axiom:sheaf-gluability}] Suppose \(\{s_{i} \mid s_{i} \in \underline{S}(U_{i})\}_{i \in \mathcal{I}}\) is a collection of sections that agree on pairwise intersection.

          Define \(s: U \to S\) as follows: for each \(x \in U\) choose \(i \in \mathcal{I}\) such that \(x \in U_{i}\) and set \(s(x) = s_{i}(x)\).
          Note that \(s\) is well-defined since if \(x \in U_{i} \cap U_{j}\) by hypothesis \(s_{i}(x) = s_{j}(x)\).
          We only need to show that \(s \in \underline{S}\).

          For every \(x \in U\) choose \(U_{i}\) such that \(x \in U_{i}\). Since \(s_{i}\) is locally constant there exists an open neighbourhood \(V \subseteq U_{i}\) of \(x\) such that \(s\) is constant in \(V\). Then since \(s|_{V} = s_{i}|_{V}\) we have that \(s\) is constant in \(V\), meaning that \(s\) is locally constant in \(U\).
  \end{itemize}
\end{example}
\SetPluralName{ex:constant-sheaf}{Constant sheaves}

We can also construct \namerefs{def:sheaf} from other data:
\begin{example} \label{ex:set-product-sheaf}
  Let \(X\) be a topological space and let \(\{A_{p}\}_{p \in X}\) be a family of sets.

  Given an open set \(U \subseteq X\) consider \(\prod(U) = \prod_{p \in U} A_{p}\) and if \(V \subseteq U\) define \(\res_{U, V} (a_{p})_{p \in U} = (a_{p})_{p \in V}\), making \(\prod\) into a \nameref{def:presheaf}.

  To check that it is indeed a \nameref{def:sheaf}, consider an open covering \(\{U_{i}\}_{i \in \mathcal{I}}\) of \(U\).
  Let \(\{s_{i} \in \prod(U_{i})\}_{i \in \mathcal{I}}\) be such that for every \(i, j \in \mathcal{I} \quad s_{i}|_{U_{i} \cap U_{j}} = s_{j}|_{U_{i} \cap U_{j}}\).
  If we write \(s_{i} = (a_{i, p})_{p \in U_{i}}\) then the compatibility condition means that \(a_{i, p} = a_{j, p}\) for every \(p \in U_{i} \cap U_{j}\).
  Therefore there exists \(s = (a_{p})_{p \in U} \in \prod(U)\) such that for every \(p \in U_{i} \quad a_{p} = a_{i, p}\).
\end{example}

\begin{definition}[Pushforward \nameref{def:sheaf}] \label{def:pushforward-sheaf}
  Let \(f: X \to Y\) be a continuous map between two topological spaces \(X\) and \(Y\).
  For any \nameref{def:sheaf} \(\sheaf{F}\) over \(X\), define the \textbf{pushforward sheaf} \(f_{\ast}\sheaf{F}\) in \(Y\) as, for every open subset \(V \subseteq Y\), \(f_{\ast}\sheaf{F}(V) = \sheaf{F}(f^{-1}(V))\).

  This defines a \textit{pushforward \nameref{def:functor}} \(f_{\ast}: \Sh(X) \to \Sh(Y)\).
\end{definition}
\SetPluralName{def:pushforward-sheaf}{Pushforwards sheaves}

\begin{example}[Skyscrapper sheaf] \label{ex:skyscrapper-sheaf}
  Let \(X\) be a topological space with \(p \in X\).
  Consider the inclussion \(\iota_{p}: \{p\} \hookrightarrow X\).
  Then the \textbf{skyscrapper sheaf} \(\iota_{p, \ast} \underline{S}\) is the \nameref{def:pushforward-sheaf} of the \nameref{ex:constant-sheaf} \(\underline{S}\):
  \begin{equation*}
    \iota_{p, \ast}\underline{S}(U) = \begin{cases}
      S & p \in U \\
      \{\ast\} & p \notin U
    \end{cases}
  \end{equation*}
\end{example}
\SetPluralName{ex:skyscrapper-sheaf}{Skyscrapper sheaves}

We can extend \autoref{def:sheaf} to more general categories using the following\change{Motivation for this definition. Why is it equivalent?}
\begin{definition}[Sheaf]
  Let \(X\) be a topological space and let \(\mathsf{C}\) be a \nameref{def:category} with \namerefs{def:categorical-product}.
  A \(\mathsf{C}\)-valued \nameref{def:presheaf} \(\sheaf{F}\) is a \textbf{sheaf} if for every open set \(U\) and open covering \(\{U_{i}\}_{i \in \mathcal{I}}\) of \(U\) the diagram
  \begin{equation} \label{diagram:sheaf}
    \sheaf{F}(U) \rightarrow \prod_{i \in \mathcal{I}} \sheaf{F}(U_{i}) \rightrightarrows \prod_{(i, j) \in \mathcal{I} \times \mathcal{I}} \sheaf{F}(U_{i} \cap U_{j})
  \end{equation}
  is an \nameref{def:equalizer}.
\end{definition}

For sufficienly nice categories we have the following \cite[\href{https://stacks.math.columbia.edu/tag/0073}{Tag 0073}]{stacks-project}\unsure{Maybe use as definition?},
\begin{proposition} \label{prop:sheaves-and-functors}
  Let \(\mathsf{C}\) be a \nameref{def:category} and \(F: \mathsf{C} \to \Set\) be a \nameref{def:functor} verifying:
  \begin{enumerate}
    \item \(F\) is a \nameref{def:faithful-functor}.
    \item \(\mathsf{C}\) has \namerefs{def:categorical-limit} and \(F\) commutes with them.
    \item \(F\) \nameref{def:reflects} \namerefs{def:categorical-isomorphism}.
  \end{enumerate}
  Then a \(\mathsf{C}\)-valued \nameref{def:presheaf} \(\sheaf{F}\) is a \nameref{def:sheaf} if and only if the underlying \nameref{def:presheaf} of sets \(F(\sheaf{F})\) is a \nameref{def:sheaf}.
\end{proposition}
\begin{proof}~
  \begin{itemize}
    \item[\(\implies\)] Suppose \(\sheaf{F}\) is a \(\mathsf{C}\)-valued \nameref{def:sheaf}.

          Then, for every open set \(U\), \(\sheaf{F}(U)\) is the \nameref{def:equalizer} of the diagram \eqref{diagram:sheaf}.
          Therefore since \(F\) commutes with \namerefs{def:categorical-limit} \(F(\sheaf{F}(U))\) is the \nameref{def:equalizer} of the corresponding diagram in \(\Set\).
    \item[\(\impliedby\)] Suppose now that \(F(\sheaf{F})\) is a \nameref{def:sheaf}.

          Let \(E \in \mathsf{C}\) be the \nameref{def:equalizer} of the diagram \eqref{diagram:sheaf}.
          Since \(\sheaf{F}(U)\) is a \nameref{def:presheaf} there exists a morphism \(\sheaf{F}(U) \to E\), therefore by applying \(F\) we obtain a morphism of sets \(F(\sheaf{F}(U)) \to F(E)\).

          On the one hand since \(F\) commutes with \nameref{def:categorical-limit} \(F(E)\) is the \nameref{def:equalizer} of the corresponding diagram in \(\Set\).
          On the other since \(F(\sheaf{F})\) is a \nameref{def:sheaf}, it is the \nameref{def:equalizer} of the same diagram.
          Therefore the morphism \(F(\sheaf{F}(U)) \to  F(E)\) is an \nameref{def:categorical-isomorphism}.

          Finally since \(F\) \nameref{def:reflects} \namerefs{def:categorical-isomorphism} we have that \(\sheaf{F}(U) \to E\) is an \nameref{def:categorical-isomorphism}.
  \end{itemize}
\end{proof}
In particular \autoref{prop:sheaves-and-functors} applies to \namerefs{def:sheaf} of groups, rings, algebras/modules/vector spaces over a fixed ring, etc.

Since we will be mainly interested in \namerefs{def:sheaf} of abelian groups and rings, we can restrict ourselves to the study of \namerefs{def:sheaf} of sets and apply \autoref{prop:sheaves-and-functors} when needed.



\section{Morphisms of (pre)sheaves}

\begin{definition}[Morphism of \namerefs{def:presheaf}] \label{def:presheaf-morphism}
  Let \(\mathcal{F}, \mathcal{G}\) be two \(\mathsf{C}\)-valued \namerefs{def:presheaf} over a topological space \(X\).
  A \textbf{presheaf morphism} \(\phi: \mathcal{F} \to \mathcal{G}\) is a \nameref{def:natural-transformation} between the \namerefs{def:functor} \(\mathcal{F}\) and \(\mathcal{G}\).

  Less abstractly, a presheaf morphism \(\phi: \mathcal{F} \to \mathcal{G}\) is the a collection of maps \(\{\phi(U): \mathcal{F}(U) \to \mathcal{G}(U)\}_{U \in \open(X)}\) such that if \(V \subseteq U\) then the following diagram commutes
  \begin{center}
     \begin{tikzcd}
      \mathcal{F}(U) \arrow[r, "\phi(U)"] \arrow[d, "\res_{U, V}"] & \mathcal{G}(U) \arrow[d, "\res_{U, V}"] \\
      \mathcal{F}(V) \arrow[r, "\phi(V)"] & \mathcal{G}(V)
    \end{tikzcd}
  \end{center}
\end{definition}
\SetPluralName{def:presheaf-morphism}{Morphisms of \namerefs{def:presheaf}}


\begin{definition}[Morphism of \namerefs{def:sheaf}] \label{def:sheaf-morphism}
  Let \(\sheaf{F}, \sheaf{G}\) be two \(\mathsf{C}\)-valued \namerefs{def:sheaf} over a topological space \(X\).
  A \textbf{sheaf morphism} \(\phi: \sheaf{F} \to \sheaf{G}\) is a \nameref{def:presheaf-morphism} between \(\sheaf{F}\) and \(\sheaf{G}\).
\end{definition}
\SetPluralName{def:sheaf-morphism}{Morphisms of \nameref{def:sheaf}}

With this choice of morphisms, \namerefs{def:presheaf} and \namerefs{def:sheaf} valued in a \nameref{def:category} \(\mathsf{C}\) over a topological space \(X\) form a \nameref{def:category}.

\begin{notation}
  We will denote the \nameref{def:category} of \(\mathsf{C}\)-valued \namerefs{def:presheaf} over a topological space \(X\) as \(\PSh_{\mathsf{C}}(X)\) and the \nameref{def:category} of \(\mathsf{C}\)-valued \namerefs{def:sheaf} over \(X\) as \(\Sh_{\mathsf{C}}(X)\).
\end{notation}

\begin{remark} \label{remark:presheaf-full-subcategory}
  The category \(\Sh_{\mathsf{C}}(X)\) is a \nameref{def:full-subcategory} of \(\PSh_{\mathsf{C}}(X)\) by the definition of \nameref{def:sheaf-morphism}, i.e: the inclusion functor \(\iota: \Sh_{\mathsf{C}}(X) \hookrightarrow \PSh_{\mathsf{C}}(X)\) is a \nameref{def:fully-faithful-functor}.
\end{remark}

\begin{definition}[\(\HomS\) \nameref{def:sheaf}] \label{def:hom-sheaf}
  Let \(\sheaf{F}, \sheaf{G} \in \Sh_{\Set}(X)\) be two \namerefs{def:sheaf} of sets over \(X\).
  We define the \(\HomS\) \nameref{def:sheaf} (of sets) between \(\sheaf{F}\) and \(\sheaf{G}\) as
  \begin{equation*}
    \forall U \in \open(X) \quad \HomS(\sheaf{F}, \sheaf{G})(U) = \Mor(\sheaf{F}|_{U}, \sheaf{G}|_{U})
  \end{equation*}
\end{definition}
\SetPluralName{def:hom-sheaf}{\(\HomS\) \namerefs{def:sheaf}}

\begin{remark}
  Let's check that the \nameref{def:hom-sheaf} \(\HomS(\sheaf{F}, \sheaf{G})\) is indeed a \nameref{def:sheaf} of sets for any \(\sheaf{F}, \sheaf{G} \in \Sh_{\Set}(X)\).

  First, we need to verify that it is a \nameref{def:presheaf}.

  We only need to define the restriction maps and check commutativity.
  For a fixed open set \(U\), the restriction of a section \(s \in \HomS(\sheaf{F}, \sheaf{G})(U)\) to an open subset \(V \subseteq U\) is \(s|_{V} \in \Mor(\sheaf{F}|_{V}, \sheaf{G}|_{V})\) defined as \(s|_{V}(W) = s(W): \sheaf{F}(W) \to \sheaf{G}(W)\) for every open set \(W \subseteq V\).
  Note that, since \(W\) is an open subset of \(U\), commutativity in \(V\) is inherited from commutativity in \(U\).

  Now, we will check the \nameref{def:sheaf} axioms. To do this, consider an open covering \(\{U_{i}\}_{i \in \mathcal{I}}\) of an open set \(U\).
  \begin{itemize}
    \item[\ref{axiom:sheaf-identity}] Let \(\phi, \psi \in \HomS(\sheaf{F}, \sheaf{G})(U)\) be two sections such that for every \(i \in \mathcal{I} \;\; \phi|_{U_{i}} = \psi|_{U_{i}}\). We want to show that for every open subset \(V \subseteq U\) the maps \(\phi(V)\) and \(\psi(V)\) are the same.

          First, note that \(\{V \cap U_{i}\}_{i \in \mathcal{I}}\) is an open cover of \(V\).
          Then, for any \(s \in \sheaf{F}(V)\) consider the restrictions \(s|_{V \cap U_{i}} \in \sheaf{F}(V \cap U_{i})\).

          Since \(\phi|_{U_{i}} = \psi|_{U_{i}}\) we have that
          \begin{equation*}
            \phi(V)(s|_{V \cap U_{i}}) = \phi(V \cap U_{i})(s|_{V \cap U_{i}}) = \psi(V \cap U_{i})(s|_{V \cap U_{i}}) = \psi(V)(s|_{V \cap U_{i}})
          \end{equation*}

          Therefore, by applying \ref{axiom:sheaf-identity} to \(\sheaf{G}\) we have that \(\phi(V)(s) = \psi(V)(s)\), i.e: \(\phi(V) = \psi(V)\).
    \item[\ref{axiom:sheaf-gluability}] Let \(\{\phi_{i} \mid \phi_{i} \in \HomS(\sheaf{F}, \sheaf{G})(U_{i})\}\) be a collection of sections that agree on pairwise intersection. We want to show that there exists \(\phi \in \HomS(\sheaf{F}, \sheaf{G})(U)\) such that for every \(i \in \mathcal{I} \;\; \phi|_{U_{i}} = \phi_{i}\).

          Let \(V \subseteq U\) be an open set and \(s \in \sheaf{F}(V)\).
          For every \(i \in \mathcal{I}\) consider \(s_{i} = s|_{V \cap U_{i}} \in \sheaf{F}(V \cap U_{i})\).
          Applying \(\phi_{i}\) we get \(t_{i} = \phi_{i}(V \cap U_{i})(s_{i}) \in \sheaf{G}(V \cap U_{i})\).

          Since for any \(i, j \in \mathcal{I} \;\; \phi_{i}|_{U_{i} \cap U_{j}} = \phi_{j}|_{U_{i} \cap U_{j}}\) we have that \(t_{i}|_{V \cap U_{i} \cap U_{j}} = t_{j}|_{V \cap U_{i} \cap U_{j}}\).
          Therefore since the sections \(\{t_{i}\}_{i \in \mathcal{I}}\) agree on pairwise intersection applying \ref{axiom:sheaf-gluability} we have that there exists \(t \in \sheaf{G}(V)\) such that \(\forall i \in \mathcal{I} \;\; t|_{V \cap U_{i}} = t_{i}\).

          This defines a map \(\phi(V): \sheaf{F}(V) \to \sheaf{G}(V)\) for every open set \(V \subseteq U\).
          We need to show that this map is a \nameref{def:sheaf-morphism}, i.e: that it commutes with restriction maps.
          Let \(W \subseteq V\) be an open set and \(s \in \sheaf{F}(V)\). We want to show that \(\phi(W)(s|_{W}) = \phi(V)(s)|_{W}\).

          By construction, \(\phi(V)(s)\) is the unique section in \(\sheaf{G}(V)\) such that for all \(i \in \mathcal{I}\)
          \begin{equation*}
            \phi(V)(s)|_{V \cap U_{i}} = \phi(V \cap U_{i})(s|_{V \cap U_{i}})
          \end{equation*}
          Then restricting boths sides to \(W \cap U_{i}\)
          \begin{equation*}
            \left( \phi(V \cap U_{i})(s|_{V \cap U_{i}}) \right)|_{W \cap U_{i}} = \left( \phi(V \cap U_{i})(s|_{V \cap U_{i}}) \right)|_{W \cap U_{i}} = \phi_{i}(V \cap U_{i})(s|_{V \cap U_{i}})|_{W \cap U_{i}}
          \end{equation*}
          Since \(\phi_{i}\) is a \nameref{def:sheaf-morphism} it commutes with restrictions, giving
          \begin{equation*}
             \phi_{i}(V \cap U_{i})(s|_{V \cap U_{i}})|_{W \cap U_{i}} = \phi_{i}(W \cap U_{i})((s|_{V \cap U_{i}})|_{W \cap U_{i}}) = \phi_{i}(W \cap U_{i})(s|_{W \cap U_{i}})
           \end{equation*}
          Therefore
          \begin{equation*}
            \phi(V)(s)|_{W \cap U_{i}} = \phi_{i}(W \cap U_{i})(s|_{W \cap U_{i}})
          \end{equation*}

          On the other hand, \(\phi(W)(s|_{W})\) is by definition the unique section in \(\sheaf{G}(W)\) such that for all \(i \in \mathcal{I}\)
          \begin{equation*}
            \phi(W)(s|_{W})|_{W \cap U_{i}} = \phi_{i}(W \cap U_{i})((s|_{W})|_{W \cap U_{i}}) = \phi_{i}(W \cap U_{i})(s|_{W \cap U_{i}})
          \end{equation*}
          Therefore we have that for all \(i \in \mathcal{I}\)
          \begin{equation*}
            \phi(V)(s)|_{W \cap U_{i}} = \phi(W)(s|_{W})|_{W \cap U_{i}}
          \end{equation*}
          Applying \ref{axiom:sheaf-identity} we get \(\phi(V)(s)|_{W} = \phi(W)(s|_{W})\).

          It only remains to show that \(\phi|_{U_{i}} = \phi_{i}\).
          Indeed, for every open subset \(V \subseteq U_{i}\) and \(s \in \sheaf{F}(V)\) we have that \(\phi(V)(s) = \phi_{i}(V)(s)\) so applying \ref{axiom:sheaf-gluability} on the open cover \(\{V \cap U_{j}\}_{j \in \mathcal{I}}\) we have that \(\phi(V) = \phi_{i}(V)\).
  \end{itemize}

  Note that to check the \nameref{def:sheaf} axioms we only used that \(\sheaf{G}\) is a \nameref{def:sheaf}, therefore the hypothesis of \(\sheaf{F}\) being a \nameref{def:sheaf} may be dropped.
\end{remark}

\begin{remark}\improvement{This should be rewritten.}
  We can extend \nameref{def:hom-sheaf} to other evaluation \namerefs{def:category} other than \(\Set\).
  However, this will be a sheaf only if the evaluation \nameref{def:category} is \textit{sufficiently} nice (in particular, if it is complete).

  For example, if \(\sheaf{F}, \sheaf{G} \in \Sh_{\Ab}(X)\) we can define the \(\HomS\) sheaf, this time an abelian group \nameref{def:sheaf}, by considering the maps \(\sheaf{F}|_{U} \to \sheaf{G}|_{U}\) as morphisms between abelian groups.
  To distinguish these \namerefs{def:hom-sheaf} from the \(\Set\)-valued version, we will denote them by \(\HomS_{\Sh_{\mathsf{C}}(X)}\), although many times the subscript will be dropped if the category is clear.
\end{remark}

\begin{example}
  If \(\sheaf{F}\) is a \nameref{def:sheaf} of sets over a topological space \(X\) then the \nameref{def:hom-sheaf} is \(\HomS(\underline{\{p\}}, \sheaf{F}) \cong \sheaf{F}\).

  First, note that for every open set \(U \subseteq X \;\; \HomS(\underline{\{p\}}, \sheaf{F})(U) = \Hom(\underline{\{p\}}|_{U}, \sheaf{F}|_{U})\).
  Therefore for every \(U \subseteq X\) the elements of the \nameref{def:hom-sheaf} at \(U\) are morphisms\\ \(\phi: \underline{\{p\}}|_{U} \to \sheaf{F}|_{U}\), i.e: a collection of maps \(\phi(V): \underline{\{p\}}|_{V} \to \sheaf{F}|_{V}\) that commute with the restriction maps for every \(V \subseteq U\) open.
  Furthermore since \(\underline{\{p\}}|_{V} = \{p\}\) the maps \(\phi(V)\) are determined by the image of \(p\).

  Consider \(s^{(V)} = \phi(V)(p) \in \sheaf{F}(V)\).
  Since \(\phi\) commutes with restriction maps we have that for every \(W \subseteq V\) open \(\res_{V, W} s^{(W)} = s^{(V)}\).
  Therefore \(\{s^{(V)}\}_{V \subseteq U}\) is a compatible family of sections, which means that by \ref{axiom:sheaf-gluability} there exists a unique section \(s \in \sheaf{F}(U)\) such that \(s|_{V} = s^{(V)}\) for every \(V \subseteq U\) open.
  In particular taking \(V = U\) we have that \(s|_{U} = s\), i.e: the morphism \(\phi\) is uniquely determined by the section \(s = \phi(U)(p) \in \sheaf{F}(U)\).

  On the other hand, given any section \(s \in \sheaf{F}(U)\) we can define a morphism\\ \(\phi: \underline{\{p\}}|_{U} \to \sheaf{F}|_{U}\) by \(\phi(V)(p) = s|_{V}\) for all \(V \subseteq U\) open.

  Since this morphism satisfies the compatibility conditions, for every open set \(U \subseteq X\) we have a bijection
  \begin{align*}
    \Theta_{U}:& \HomS(\underline{\{p\}}, \sheaf{F})(U) \to \sheaf{F}(U) \\
    & \phi \mapsto \phi(U)(p)
  \end{align*}

  It remains to show that \(\Theta_{U}\) is a \nameref{def:sheaf-morphism}, i.e: that the following diagram commutes for every \(V \subseteq U\) open.
  \begin{center}
   \begin{tikzcd}
     \HomS(\underline{\{p\}}, \sheaf{F})(U) \arrow[r, "\Theta_{U}"] \arrow[d, "\res_{U, V}"] & \sheaf{F}(U) \arrow[d, "\res_{U, V}"] \\
     \HomS(\underline{\{p\}}, \sheaf{F})(V) \arrow[r, "\Theta_{V}"] & \sheaf{F}(V)
   \end{tikzcd}
  \end{center}

  Let \(\phi \in \HomS(\underline{\{p\}}, \sheaf{F})(U)\).
  Note that since \(\phi\) is a \nameref{def:sheaf-morphism}
  \begin{align*}
    \left( \Theta_{V} \circ \res_{U, V} \right)(\phi) &= \Theta_{V}(\phi|_{V}) = (\phi|_{V})(V)(p) = \phi(V)(p) \\
                                          &= \res_{U, V}\phi(U)(p) = (\res_{U, V} \circ \Theta_{U})(\phi)
  \end{align*}

  Therefore the collection of maps \(\{\Theta_{U}\}_{U \in \open(X)}\) defines an \nameref{def:categorical-isomorphism}
  \begin{equation*}
    \HomS(\underline{\{p\}}, \sheaf{F}) \cong \sheaf{F}
  \end{equation*}
\end{example}

\begin{example}
  If \(\sheaf{F}\) is a \nameref{def:sheaf} of abelian groups then \(\HomS(\underline{\mathbb{Z}}, \sheaf{F}) \cong \sheaf{F}\).
\end{example}

\section{Stalks and germs.}

\begin{definition}[Stalk] \label{def:stalk}
  Let \(\mathcal{F}\) be a \nameref{def:presheaf} of sets over a topological space \(X\).
  The \textbf{stalk} of \(\mathcal{F}\) at a point \(p \in X\) is
  \begin{equation*}
    \mathcal{F}_{p} = \colim_{\substack{U \ni p \\ U \text{ open}}} \mathcal{F}(U)
  \end{equation*}
  where the \nameref{def:categorical-colimit} is over the set of open neighbourhoods of \(p\).

  This defines a \textit{stalkification \nameref{def:functor}} \((-)_{p}: \PSh_{\Set}(X) \to \Set\).
\end{definition}
\begin{remark} \label{remark:stalk}
  Since the set of open neighbourhoods of a point is partially ordered by reverse inclussion the \nameref{def:stalk} \nameref{def:categorical-colimit} is a \nameref{def:directed-colimit}.
  Therefore we can describe \(\mathcal{F}_{p}\) as
  \begin{equation*}
    \mathcal{F}_{p} = \quot{\{(s, U) \mid p \in U, \, s \in \mathcal{F}(U)\}}{\sim}
  \end{equation*}
  where \((s, U) \sim (t, V)\) if and only if there exists an open set \(W \subseteq U \cap V\) with \(p \in W\) such that \(s|_{W} = t|_{W}\).

  For a section \(f \in \mathcal{F}(U)\) we call \(f_{p} = [(f, U)] \in \mathcal{F}_{p}\) the \textbf{germ} of \(f\) at \(p\).
\end{remark}

\begin{example}
  Following \autoref{ex:set-product-sheaf} given a \nameref{def:presheaf} of sets over a topological space \(X\) we can define the \nameref{def:sheaf} \(\prod(\mathcal{F})\) as \(\prod(\mathcal{F})(U) = \prod_{p \in U} \mathcal{F}_{p}\) for every open set \(U \subseteq X\).
\end{example}


\begin{proposition} \label{prop:map-to-prod-stalks-injective}
  Let \(X\) be a topological space and let \(\sheaf{F}\) be a \nameref{def:sheaf} of sets over \(X\). Then for every open set \(U \subseteq X\) the map \(\sheaf{F}(U) \to \prod_{p \in U} \sheaf{F}_{p}\) is injective.
\end{proposition}
\begin{proof}
  Let \(s, t \in \sheaf{F}(U)\) be two sections that map to the same element in every \nameref{def:stalk} \(\sheaf{F}_{p}\) for every \(p \in U\).
  Then by \autoref{remark:stalk} for every \(p \in U\) there exists an open neighbourhood \(V^{p} \subseteq U\) such that \(s|_{V^{p}} = t|_{V^{p}}\).
  Therefore since \(\{V_{p}\}_{p \in U}\) is an open covering of \(U\) by \ref{axiom:sheaf-identity} \(s = t\).
\end{proof}

\begin{proposition}[Induced \nameref{def:stalk} morphism] \label{prop:induced-stalk-morphism}
  Let \(\mathcal{F}, \mathcal{G} \in \PSh_{\Set}(X)\) be two \namerefs{def:presheaf} of sets over \(X\) and \(\phi: \mathcal{F} \to \mathcal{G}\) a \nameref{def:presheaf-morphism}.
  Then for every \(p \in X\), \(\phi\) induces a morphism of \namerefs{def:stalk} \(\phi_{p}: \mathcal{F}_{p} \to \mathcal{G}_{p}\).
\end{proposition}
\begin{proof}
  We have the following diagram, where the vertical arrows come from the \nameref{def:categorical-colimit} definition.

  \begin{center}
    \begin{tikzcd}
      \mathcal{F}(U) \arrow[r, "\phi(U)"] \arrow[d] & \mathcal{G}(U) \arrow[d] \\
      \mathcal{F}_{p} \arrow[r, "\phi_{p}", dashed] & \mathcal{G}_{p}
    \end{tikzcd}
  \end{center}
  Therefore, we just need to define for every \([(f, U)] \in \mathcal{F}_{p} \quad \phi_{p}([(f, U)]) = [(\phi(U)(f), U)]\).
\end{proof}


To generalize the concept of \nameref{def:stalk} to other \namerefs{def:category} in a similar fashion to \autoref{prop:sheaves-and-functors} we have the following result
\begin{proposition}
  Let \(\mathsf{C}\) be a \nameref{def:category} and \(F: \mathsf{C} \to \Set\) be a \nameref{def:functor} verifying:
  \begin{enumerate}
    \item \(F\) is a \nameref{def:faithful-functor}.
    \item \namerefs{def:directed-colimit} exist in \(\mathsf{C}\) and \(F\) commutes with them.
  \end{enumerate}
  Let \(X\) be a topological space and \(p \in X\). Let \(\mathcal{F}\) be a \(\mathsf{C}\)-valued \nameref{def:presheaf}.
  Then \(\mathcal{F}_{p} = \colim_{p \in U} \mathcal{F}(U)\) exists in \(\mathsf{C}\) and its underlying set \(F(\mathcal{F}_{p})\) is equal to the \nameref{def:stalk} of sets of \(\mathcal{F}\).

  Furthermore the construction \(\mathcal{F} \mapsto \mathcal{F}_{p}\) is functorial.
\end{proposition}
\begin{proof}
  Since \namerefs{def:directed-colimit} exist in \(\mathsf{C}\) by \autoref{remark:stalk} the \nameref{def:stalk} \(\mathcal{F}_{p}\) exists in \(\mathsf{C}\).
  Uniqueness of the \nameref{def:stalk} follows analogously to \autoref{prop:sheaves-and-functors}.
\end{proof}

\begin{theorem}[\namerefs{def:stalk} and \namerefs{def:categorical-isomorphism}] \label{thm:induced-stalk-isomorphism}
  Let \(\sheaf{F}, \sheaf{G}\) be two \namerefs{def:sheaf} of sets over a topological space \(X\) and let \(\phi: \sheaf{F} \to \sheaf{G}\) be a \nameref{def:sheaf-morphism}.
  Then \(\phi\) is an \nameref{def:categorical-isomorphism} if and only if for every \(p \in X\) the \nameref{prop:induced-stalk-morphism} \(\phi_{p}: \sheaf{F}_{p} \to \sheaf{G}_{p}\) is an \nameref{def:categorical-isomorphism}.
\end{theorem}
\begin{proof}
  Suppose first that \(\phi\) is an \nameref{def:categorical-isomorphism}.
  Then since taking \namerefs{def:stalk} is a \nameref{def:functor} and \nameref{prop:functors-preserve-isomorphisms}, \(\phi_{p}\) is an \nameref{def:categorical-isomorphism}.

  To prove the converse, suppose that for every \(p \in X \;\; \phi_{p}\) is an \nameref{def:categorical-isomorphism}.
  We are going to show that for every open set \(U \subseteq X \;\; \phi(U): \sheaf{F}(U) \to \sheaf{G}(U)\) is an \nameref{def:categorical-isomorphism} by proving that it is an injective and surjective morphism in \(\Set\).

  For the injectivity, let \(s, t \in \sheaf{F}(U)\) be two sections such that \(\phi(U)(s) = \phi(U)(t)\).
  Then for every \(p \in U\) we have that \(\phi_{p}(s_{p}) = \phi_{p}(t_{p})\).
  Since \(\phi_{p}\) is an \nameref{def:categorical-isomorphism} in particular is injective, therefore \(s_{p} = t_{p}\).
  This means that there exist \(V_{p} \subseteq U\) open neighbourhood of \(p\) such that \(s|_{V_{p}} = t|_{V_{p}}\).
  Since \(\{V_{p}\}_{p \in U}\) is an open cover of \(U\) by \ref{axiom:sheaf-identity} we have that \(s = p\), i.e: \(\phi(U)\) is injective.

  For surjectivity, consider \(t \in \sheaf{G}(U)\).
  For every \(p \in U\) since \(\phi_{p}\) is surjective we have that there exists \(s_{p} \in \sheaf{F}_{p}\) such that \(\phi_{p}(s_{p}) = t_{p}\).
  By definition of \nameref{def:stalk} we have that there exists an open neighbourhood \(U_{p} \subseteq U\) of \(p\) and a section \(s^{(p)} \in \sheaf{F}(U_{p})\) such that \(\phi(U_{p})(s^{(p)}) = t|_{U_{p}}\).
  Note that on overlaps \(U_{p} \cap U_{q}\) we have that
  \begin{equation*}
    \phi(U_{p} \cap U_{q})(s^{(p)}|_{U_{p} \cap U_{q}}) = t|_{U_{p} \cap U_{q}} = \phi(U_{p} \cap U_{q})(s^{(q)}|_{U_{p} \cap U_{q}})
  \end{equation*}
  Therefore since \(\phi\) is injective we have that \(s^{(p)}|_{U_{p} \cap U_{q}} = s^{(q)}|_{U_{p} \cap U_{q}}\).
  Then by \ref{axiom:sheaf-gluability} there exists a section \(s \in \sheaf{F}(U)\) such that \(\phi(U)(s) = t\), proving surjectivity.
\end{proof}


\section{Sheafification}
In \autoref{ex:constant-presheaf} we have seen how we can \textit{convert} a \nameref{def:presheaf} into a \nameref{def:sheaf}.
We are now going to see how we can construct \namerefs{def:sheaf} from \namerefs{def:presheaf}.

We will first define categorically what \textbf{sheafification} ought to be and then we will construct it, therefore showing its existence.

We can restrict ourselves to the case where \(\mathsf{C} = \Set\) and use \autoref{prop:sheaves-and-functors} when working on more \textit{general} categories.

\begin{definition}[Sheafification] \label{def:sheafification}
  Let \(\mathcal{F}\) be a \nameref{def:presheaf} of sets over a topological space \(X\).
  A \nameref{def:presheaf-morphism} \((-)^{\#}: \mathcal{F} \to \mathcal{F}^{\#}\) is the \textbf{sheafification} of \(\mathcal{F}\) if \(\mathcal{F}^{\#}\) is a \nameref{def:sheaf} and for any other \nameref{def:sheaf} \(\sheaf{G}\) and \nameref{def:presheaf-morphism} \(\psi: \mathcal{F} \to \sheaf{G}\) there exists a unique \nameref{def:sheaf-morphism} \(\phi: \mathcal{F}^{\#} \to \sheaf{G}\) making the following diagram commute
  \begin{center}
    \begin{tikzcd}
      \mathcal{F} \arrow[r, "(-)^{\#}"] \arrow[rd, "\psi" below] & \mathcal{F}^{\#} \arrow[d, "\phi"] \\
      & \sheaf{G}
    \end{tikzcd}
  \end{center}
  This defines a \textit{sheafification \nameref{def:functor}} \((-)^{\#}: \PSh_{\Set}(X) \to \Sh_{\Set}(X)\).
\end{definition}

\begin{remark}[Uniqueness of \nameref{def:sheafification}]
  Since \nameref{def:sheafification} is defined using an \nameref{def:universal-property} if it exists it is unique (up to unique \nameref{def:categorical-isomorphism}).
  Therefore we can speak of \textit{the} \nameref{def:sheafification} of a \nameref{def:presheaf} instead of \textit{a} \nameref{def:sheafification}.
\end{remark}

\begin{remark} \label{remark:sheafification-inclussion-adjunction}
  The \nameref{def:sheafification} functor is left-\nameref{def:adjoint} to the inclussion functor \(i: \Sh_{\Set}(X) \hookrightarrow \PSh_{\Set}(X)\).
  This means that for every \nameref{def:presheaf} \(\mathcal{F}\) and \nameref{def:sheaf} \(\sheaf{G}\) we have \(\Hom_{\Sh}(\mathcal{F}^{\#}, \sheaf{G}) \cong \Hom_{\PSh}(\mathcal{F}, i(\sheaf{G}))\).
\end{remark}


\begin{example}
  For any \nameref{def:sheaf} \(\sheaf{F}\) its \nameref{def:sheafification} is itself: \(\sheaf{F}^{\#} = \sheaf{F}\) with \((-)^{\#}: \sheaf{F} \to \sheaf{F}^{\#}\) the identity morphism.
\end{example}

Before constructing the \nameref{def:sheafification} functor we will take a detour to study the inverse image \nameref{def:sheaf} and its relation to \namerefs{def:stalk}.

\begin{lemma}
  Let \(f: X \to Y\) be a continuous map between two topological spaces \(X\) and \(Y\).
  Let \(\sheaf{G}\) be a \nameref{def:sheaf} over \(Y\).

  Then \(f_{\pre}^{-1}\sheaf{G}\) defined as for every open set \(U \subseteq X \quad f_{\pre}^{-1}\sheaf{G}(U) = \colim_{V \supset f(U)} \sheaf{G}(V)\) is a \nameref{def:presheaf} but not a \nameref{def:sheaf}.
\end{lemma}
\begin{proof}
\end{proof}


\begin{definition}[Inverse image \nameref{def:sheaf}] \label{def:inverse-image-sheaf}
  Let \(f: X \to Y\) be a continuous map between two topological spaces \(X\) and \(Y\).
  Let \(\sheaf{G}\) be a \nameref{def:sheaf} of sets over \(Y\).
  Then \textbf{inverse image} \nameref{def:sheaf} of \(\sheaf{G}\) by \(f\) is \(f^{-1}\sheaf{G} = (f_{\pre}^{-1}\sheaf{G})^{\#}\).

  This defines an \textit{inverse image \nameref{def:functor}} \(f^{-1}: \Sh(Y) \to \Sh(X)\).
\end{definition}

\begin{example}
  Let \(Y\) be a topological space and \(p \in Y\).
  Then the \nameref{def:inverse-image-sheaf} \(f^{-1}\sheaf{G}\) corresponds to the \nameref{def:stalk} \(\sheaf{G}_{p}\).
\end{example}

\begin{proposition} \label{prop:pushforward-inverse-adjunction}
  For every continuous map \(f: X \to Y\) the functors induced by the \nameref{def:pushforward-sheaf} and the \nameref{def:inverse-image-sheaf} are \nameref{def:adjoint}.
\end{proposition}

\begin{remark}
  \autoref{prop:pushforward-inverse-adjunction} tells us that taking \namerefs{def:stalk} and the \nameref{ex:skyscrapper-sheaf} are \nameref{def:adjoint} \namerefs{def:functor}.
\end{remark}


We will now construct the \nameref{def:sheafification} of a \nameref{def:presheaf} \(\mathcal{F}\).
Our goal is to prove the following.
\begin{theorem}[Construction of the \nameref{def:sheafification}]\label{thm:sheafification-construction}
  Let \(\mathcal{F}\) be a \nameref{def:presheaf} of sets over a topological space \(X\). Then
  \begin{equation} \label{eq:f-sharp}
    \mathcal{F}^{\#}(U) = \left\{ (f_{p} \in \mathcal{F}_{p})_{p \in U} \mid f_{p} \text{ verifies } (\ast) \right\}
  \end{equation}
  where \((\ast)\) is the property "for all \(p \in U\) there exists an open neighbourhood \(V \subseteq U\) of \(p\) and \(s \in \mathcal{F}(V)\) such that for all \(q \in V \quad f_{q} = s_{q}\)" is the \nameref{def:sheafification} of \(\mathcal{F}\).
\end{theorem}

We will prove this in steps in the next lemmas.

\begin{lemma}
  \(\mathcal{F}^{\#}\) is a \nameref{def:sheaf} (of sets) and there exists a \textit{canonical} \nameref{def:presheaf-morphism} \(\mathcal{F} \to \mathcal{F}^{\#}\).
\end{lemma}
\begin{proof}
  Note that given \(u \in U\) the condition \((\ast)\) depends only on the values of \(s_{q}\) for \(q\) in an open neighbourhood of \(p\) the map \(\prod_{p \in U} \mathcal{F}_{p} \to \prod_{q \in V} \mathcal{F}_{q}\) sends elements of \(\mathcal{F}^{\#}(U)\) to \(\mathcal{F}^{\#}(V)\).
  Therefore taking this maps as restriction maps we make \(\mathcal{F}^{\#}\) into a \nameref{def:presheaf}.

  Furthermore using the map described in \autoref{prop:map-to-prod-stalks-injective} for every \(V \subseteq U\) open we get the following commutative diagram
  \begin{center}
    \begin{tikzcd}
      \mathcal{F}(U) \arrow[r] \arrow[d] & \mathcal{F}^{\#}(U) \arrow[r, tail] \arrow[d] & \prod\limits_{p \in U} \mathcal{F}_{p} \arrow[d] \\
      \mathcal{F}(V) \arrow[r] & \mathcal{F}^{\#}(V) \arrow[r, tail] & \prod\limits_{q \in V} \mathcal{F}_{q}
    \end{tikzcd}
  \end{center}
  where the vertical arrows come from the restriction maps, giving a \nameref{def:presheaf-morphism} \(\mathcal{F} \to \mathcal{F}^{\#}\).

  To see that \(\mathcal{F}^{\#}\) is indeed a \nameref{def:sheaf}, consider an open set \(U \subseteq X\) and an open covering \(\{U_{i}\}_{i \in \mathcal{I}}\) of \(U\).

  For every \(i \in \mathcal{I}\) let \(s_{i} = (s_{i, p})_{p \in U_{i}} \in \mathcal{F}^{\#}(U_{i})\) such that for all \(i, j \in \mathcal{I} \quad s_{i}|_{U_{i} \cap U_{j}} = s_{j}|_{U_{i} \cap U_{j}}\).
  Since \(\prod_{p \in U} \mathcal{F}_{p}\) is a \nameref{def:sheaf} we can get an element \(s = (s_{p})_{p \in U} \in \prod_{p \in U} \mathcal{F}_{p}\) by restricting \(s_{i}\) to \(U_{i}\).
  Consider \(p \in U\). Since \(p \in U_{i}\) for some \(i \in \mathcal{I}\) by \((\ast)\) there exists an open neighbourhood \(V\) of \(p\) and \(t \in \mathcal{F}(V)\) such that for all \(q \in V \quad s_{i, q} = (t, V)\).
  Therefore since \(s_{i, q} = s_{q}\) we get that \(s\) verifies \((\ast)\).
\end{proof}

\begin{lemma} \label{lemma:stalks-commute-sheafification}
  Let \(X\) be a topological space and let \(\mathcal{F}\) be a \nameref{def:presheaf} of sets over \(X\). Then for every \(p \in X\) we have \(\mathcal{F}_{p} \cong \mathcal{F}_{p}^{\#}\).
\end{lemma}
\begin{proof}
  Let \(S\) be an arbitrary set.
  Then using the \namerefs{def:adjoint} \autoref{remark:sheafification-inclussion-adjunction} and \autoref{prop:pushforward-inverse-adjunction} we have
  \begin{align*}
    \Hom_{\Set}((\mathcal{F}^{\#})_{p}, S) &\cong \Hom_{\Sh}(\mathcal{F}^{\#}, \iota_{p, \ast}\underline{S}) \\
                                 &\cong \Hom_{\PSh}(\mathcal{F}, \iota_{p, \ast}\underline{S}) \\
                                 &\cong \Hom_{\Set}(\mathcal{F}_{p}, S)
  \end{align*}
  Therefore since this holds for every set \(S\) we have that the functors \(\Hom_{\Set}((\mathcal{F}^{\#})_{p}, -)\) and \(\Hom_{\Set}(\mathcal{F}_{p}, -)\) are isomorphic.
  Finally applying the \nameref{lemma:Yonneda} we get \((\mathcal{F}^{\#})_{p} \cong \mathcal{F}_{p}\).
\end{proof}


\begin{lemma}
  Let \(\sheaf{G}\) be a \nameref{def:sheaf} of sets and let \(\mathcal{F} \to \sheaf{G}\) be a \nameref{def:presheaf-morphism}. Then \(\mathcal{F} \to \sheaf{G}\) factors uniquely as \(\mathcal{F} \to \mathcal{F}^{\#} \to \sheaf{G}\).
\end{lemma}
\begin{proof}
  Consider the following commutative diagram
  \begin{center}
    \begin{tikzcd}
      \mathcal{F} \arrow[r] \arrow[d] & \mathcal{F}^{\#} \arrow[r] \arrow[d] & \prod(\mathcal{F}) \arrow[d] \\
      \mathcal{G} \arrow[r] & \mathcal{G}^{\#} \arrow[r] & \prod(\mathcal{G})
    \end{tikzcd}
  \end{center}
  We only need to prove that \(\sheaf{G} \cong \sheaf{G}^{\#}\).
  By \nameref{thm:induced-stalk-isomorphism} we only need to show that for every \(p \in X\) the \nameref{prop:induced-stalk-morphism} is an \nameref{def:categorical-isomorphism}.
  But this is precisely \autoref{lemma:stalks-commute-sheafification}.
\end{proof}

\unsure[inline]{Should we mention the \textbf{espace étalé}?}



\end{document}
